\title{Parameter-Efficient Orthogonal Finetuning via Factorization}

\begin{abstract}
Large foundation models have become widespread, but their full-scale training remains computationally intensive. Consequently, optimizing the adaptation of these models to specific tasks is of growing importance. In this work, we examine Orthogonal Finetuning (OFT), a systematic approach for task adaptation. Although OFT achieves strong generalization, it requires a substantial number of trainable parameters owing to the inherently high dimensionality of orthogonal matrices. To mitigate this, we analyze OFT from an information transmission standpoint and derive several core requirements for enhanced parameter-efficiency. Drawing inspiration from the efficiency of the Cooley-Tukey fast Fourier transform algorithm, we introduce a compact orthogonal parameterization leveraging butterfly structures. Applying this parameterization within the OFT paradigm, we develop Orthogonal Butterfly~(BOFT), a novel finetuning approach that emphasizes parameter efficiency. BOFT generalizes the OFT framework by treating it as a particular instance, thereby extending the possibilities of orthogonal finetuning. We further present comprehensive experimental results, demonstrating the adaptation of large vision transformers, large language models, and text-to-image diffusion models across a spectrum of tasks in vision and language.
\end{abstract}

\section{Introduction}

State-of-the-art models such as ChatGPT~\cite{brown2020language,chen2021evaluating} and Stable Diffusion~\cite{rombach2022high} exemplify the outstanding generalization capacity of large-scale foundation models. The rapid escalation in these models' performance has coincided with a significant growth in parameter counts

(\emph{e.g.}, GPT-3 comprises approximately 175B parameters). This expansion has made it increasingly impractical for many to train foundation models from the ground up. Therefore, advancing the field requires effective strategies for repurposing these existing models without repeating resource-intensive pretraining. Efficient adaptation to new tasks is thus essential. Broadly, there are three major strategies for such adaptation: \emph{model finetuning}~\cite{hulora2022,zaken2022bitfit,guo2020parameter,raffel2020exploring,qiu2023controlling,zhang2022adaptive,chavan2023one}, wherein the base architecture remains fixed while selected parameters are updated; \emph{adapter tuning}~\cite{rebuffi2017learning,houlsby2019parameter,pfeiffer2020adapterfusion,lin2020exploring,he2021towards}, which introduces new trainable modules into the network; and \emph{prompt tuning}~\cite{li2021prefix,lester2021power}, where tunable prompt embeddings are prepended to model inputs. Of these, model finetuning stands out for its conceptual simplicity, effectiveness, and the notable advantage of not introducing additional inference overhead.

Model finetuning fundamentally aims to maintain a close resemblance between the pretrained and finetuned models according to specified metrics, thereby retaining the pretrained knowledge. In practice, popular finetuning techniques typically utilize a small learning rate, an empirical method intended to keep the changes between the original and adapted models minimal. Recognizing the inherent structure in weight matrices, more systematic strategies focus on preserving the relational geometry of the network, particularly the pairwise angular relationships among neurons. This principle underpins Orthogonal Finetuning (OFT)~\cite{qiu2023controlling}, a recently proposed framework where all neurons within a single layer are mapped via the same orthogonal matrix, ensuring that the angles between neuron vectors are mathematically preserved during finetuning. OFT has shown notable success in terms of generalization and convergence when applied to text-to-image diffusion model adaptation~\cite{qiu2023controlling}. However, a key challenge is that the orthogonal matrices, due to their high dimensionality, can lead to a large number of learnable parameters. To alleviate this, OFT employs a block-diagonal orthogonal matrix structure, effectively decreasing parameter count. Nevertheless, this efficiency comes with drawbacks: the imposed block-diagonal pattern restricts flexibility, forcing the orthogonal transformations to act independently across different subspaces. Such structured sparsity may inject specific inductive biases (\emph{e.g.}, diminished expressive power and inability to represent standard linear mappings), and crucially, there is still no clear methodology for determining an optimal sparsity pattern.

A central challenge in this context lies in constructing a dense orthogonal matrix without sacrificing parameter efficiency. At first, this may appear unattainable, as a dense orthogonal matrix in $d$ dimensions typically entails $\mathcal{O}(d^2)$ parameters. However, our method innovatively addresses this by expressing a dense orthogonal matrix as a product of multiple, structured sparse matrices. The primary insight behind this strategy is that, by allowing additional computational steps, the number of tunable parameters can be considerably decreased. This is accomplished by multiplying sparse matrices during every training step to represent the orthogonal matrix. To contextualize the factorization process, we conceptualize generating a dense orthogonal matrix as analogous to transmitting information across a grid-like network. Framed this way, assembling a dense orthogonal matrix from sparse components mirrors information dissemination on a structured graph. This viewpoint facilitates the design of several forms of sparse matrix factorizations that maintain a limited set of trainable parameters, yet retain the expressive power necessary for producing dense matrices.

For parameter efficiency within this information transmission perspective, we take cues from the butterfly networks utilized in the Cooley-Tukey fast Fourier transform algorithm~\cite{cooley1965algorithm}, where rapid information flow between nodes is well established~\cite{klasing1994broadcasting}. The butterfly pattern embedded in the Cooley-Tukey scheme naturally inspires a sparse matrix decomposition, known as butterfly factorization. This allows a $d$-dimensional dense matrix to be synthesized as a product of $\mathcal{O}(\log d)$ sparse matrices, each containing only $\mathcal{O}(d)$ nonzero elements. Enforcing orthogonality within each sparse component yields an orthogonal matrix parameterization with merely $\mathcal{O}(d\log d)$ parameters—representing a dramatic reduction compared to traditional approaches. Exploiting this compact orthogonal representation, we introduce a novel and efficient finetuning approach, Orthogonal Butterfly (BOFT). BOFT generalizes the OFT method, offering a broader orthogonal finetuning framework. Both block-diagonal and butterfly architectures share an important property—sparsity—which enables a reduction in the number of adjustable parameters in orthogonal matrices. Comparing our strategy with LoRA's~\cite{hulora2022} low-rank structure reveals that low-rank and sparse matrices are two principal classes of structured matrices~\cite{candes2011robust} designed to achieve parameter efficiency.

In contrast to the block-diagonal approach adopted by OFT, which seeks a balance between expressive power and structural regularity, BOFT utilizes the butterfly structure to achieve a more continuous interpolation between matrices belonging to the full orthogonal group (ensuring expressivity) and identity matrices (ensuring regularity). This methodology facilitates the identification of a more compact subset within the orthogonal group suited for specific downstream applications. Given the established role of butterfly structures in facilitating efficient linear transforms, including the discrete Fourier and discrete cosine transforms, we posit that leveraging this structure imparts an inherent inductive bias conducive to improved generalization and reduced risk of overfitting. This insight is grounded in the observation that the butterfly structure is capable of exactly reproducing several well-known linear transforms—a feat that cannot be matched by the block-diagonal structure employed by OFT. Our main contributions are summarized as follows:

\begin{itemize}[leftmargin=*,nosep]
    \item We introduce a new perspective on parameter efficiency in orthogonal finetuning by formulating it within an information transmission framework. This reinterprets the construction of parameter-efficient, dense orthogonal matrices as an information flow problem mapped onto a grid-structured graph.
    \item Drawing inspiration from the butterfly patterns inherent to the Cooley-Tukey algorithm, we propose Orthogonal Butterfly, a novel method for orthogonal finetuning designed within this information transmission paradigm to optimize parameter efficiency.
    \item We offer theoretical analysis illuminating why BOFT achieves notable reductions in the number of learnable parameters without sacrificing expressivity or training robustness. Additionally, the matrix factorization employed by BOFT naturally supports a form of weight interpolation (refer to Figure~\ref{fig:boft_interpolation}).
    \item We pioneer the application of orthogonal finetuning techniques to a broad array of problems outside controllable text-to-image generation~\cite{qiu2023controlling}, thereby highlighting its versatility as a general-purpose finetuning strategy. In pursuit of this, BOFT is evaluated across multiple adaptation scenarios in both computer vision and natural language processing. Empirical results demonstrate that BOFT significantly outperforms current leading methods, establishing its exceptional parameter efficiency and ability to generalize.
\end{itemize}

\textbf{Parameter-efficient finetuning (PEFT)}. With the expansion and increasing capabilities of foundation models (\emph{e.g.}, \cite{brown2020language,radford2021learning,rombach2022high,kirillov2023segment}), the challenge of adapting these models to specific tasks while minimizing the number of updated parameters has drawn significant attention~\cite{treviso2023efficient,ding2023parameter,lialin2023scaling}. Within the broad landscape of PEFT approaches~\cite{houlsby2019parameter,aghajanyan2020intrinsic,hulora2022,edalati2022krona,wang2022adamix,gheini2021cross,zaken2022bitfit,guo2020parameter,sung2021training,ansell2022composable,lester2021power,li2021prefix,vu2022spot,he2021towards,mao2021unipelt,karimi2021compacter,liu2022few,sung2022lst,chen2022parameter,jia2022visual,chen2022adaptformer,zhang2022neural,jie2023fact,lian2022scaling,luo2023towards}, methods that involve reparameterizing model weights~\cite{aghajanyan2020intrinsic,hulora2022,edalati2022krona} are particularly relevant here. For instance, LoRA~\cite{hulora2022} modifies the pretrained weight matrices through the addition of a low-rank matrix product, achieving strong results for many language-based applications. Since LoRA maintains a fixed rank for the added matrices across all layers, other research~\cite{zhang2022adaptive,valipour2022dylora,zhang2023increlora} has proposed adaptive mechanisms to set the rank per layer, allowing for more efficient allocation of trainable parameters. The simplicity and effectiveness of such low-rank parameterizations have fueled their widespread adoption~\cite{dettmers2023qlora,zi2023delta,chavan2023one}. Motivated by the connection between hyperspherical energy and model generalization~\cite{liu2018learning,liu2021orthogonal}, \cite{qiu2023controlling} introduced orthogonal finetuning (OFT), which provides an alternative strategy for finetuning diffusion models for text-to-image generation. OFT operates by learning an orthogonal transformation matrix that rotates neurons within each layer, yielding both enhanced generalization and more stable optimization compared to LoRA. However, this approach typically results in a larger number of trainable parameters than LoRA. Consequently, making OFT more parameter-efficient emerges as a meaningful research objective. Further, it remains uncertain how broadly OFT can be applied to adaptation tasks outside of text-to-image diffusion modeling~\cite{qiu2023controlling}. BOFT advances OFT's parameter efficiency by leveraging butterfly matrix factorization, enabling orthogonal finetuning to be practical for a wider range of adaptation tasks.

\textbf{Butterfly structures}. The radix-2 Cooley-Tukey method systematically decomposes the $N$-point discrete Fourier transform into two transforms of length $\frac{N}{2}$ at each stage, which gives rise to a computational pattern known as the butterfly structure. This structure can be represented as a sequential product of sparse matrices—commonly referred to as a butterfly matrix. Butterfly matrices have been deployed for parameterizing orthogonal transformations to bypass pivoting in Gaussian elimination and to gain computational speed~\cite{parker1995random}, to promote stability in the training of recurrent neural architectures~\cite{jing2017tunable}, and for use in approximating kernels~\cite{munkhoeva2018quadrature}. Work by \cite{dao2019learning,dao2019kaleidoscope} demonstrates the utility of butterfly parameterizations for learning efficient linear transforms. Others~\cite{chen2022pixelated,dao2022monarch} have employed butterfly matrices for the scalable training of neural models. These structures also prove beneficial in rapid matrix-vector computation~\cite{michielssen1996multilevel,o2010algorithm}, constructing data-sparse matrix representations~\cite{li2015butterfly}, and improving network transmission~\cite{klasing1994broadcasting,li2003linear,rajasingh2016transmission}. Unlike prior efforts, our approach is dedicated to utilizing butterfly structures for the explicit purpose of improving the parameter efficiency in OFT.

\section{An Information Transmission View on Orthogonal Finetuning}

We begin by reviewing key concepts in OFT. In order to finetune pretrained weights, OFT reparameterizes the updated weight matrix as a product between a learnable orthogonal matrix and the fixed pretrained weight matrix. In contrast with LoRA, which incorporates an additive low-rank matrix for weight updates, OFT adopts a multiplicative orthogonal matrix for this purpose. While LoRA achieves parameter efficiency through a low-rank formulation, the original OFT realizes this by employing a block-diagonal orthogonal matrix~\cite{qiu2023controlling}; the parameter cost decreases as the diagonal blocks become smaller. Figure~\ref{fig:comp_lora} visually contrasts these approaches. The primary rationale behind applying an orthogonal transformation during weight matrix finetuning is to maintain the pair-wise angular relationships among neurons~\cite{liu2018learning,liu2021orthogonal,qiu2023controlling}, thereby largely retaining semantic knowledge from pretraining. Specifically, OFT involves learning an orthogonal matrix $\bm{R}\in\mathbb{R}^{d\times d}$ for a pretrained linear weight matrix $\bm{W}^0\in\mathbb{R}^{d\times n}$. This changes the forward computation from $\bm{z}=(\bm{W}^0)^\top\bm{x}$ to $\bm{z}=(\bm{R}\bm{W}^0)^\top\bm{x}$, with $\bm{x}\in\mathbb{R}^{d}$ and $\bm{z}\in\mathbb{R}^{n}$ as the input and output, respectively. For zero initialization, $\bm{R}$ is set as the identity matrix. To preserve the orthogonality of $\bm{R}$ during training, we utilize Cayley parameterization as in \cite{qiu2023controlling,liu2021orthogonal}, such that $\bm{R}=(\bm{I}+\bm{Q})(\bm{I}-\bm{Q})^{-1}$ where $\bm{Q}$ is skew-symmetric with $\bm{Q}=-\bm{Q}^\top$. To improve parameter efficiency, OFT imposes a block-diagonal form on $\bm{R}$, parameterizing it as $\text{diag}(\bm{R}_1,\bm{R}_2,\cdots,\bm{R}_r)$, where each $\bm{R}_i\in\mathcal{R}^{b\times b}$ is an independent small orthogonal matrix and $br=d$. However, while this structure reduces parameter count, it does so by assuming that a neuron's dimensions (i.e., the columns of $\bm{W}^0$) can be partitioned into $r$ groups, with each group being transformed via a distinct orthogonal matrix. Although this design has demonstrated empirical benefits, partitioning neuron dimensions simply by index is not principled, suggesting a preference for dense orthogonal matrices. This leads to the following question: \emph{Is it possible to build a dense orthogonal matrix while maintaining parameter efficiency?}

To tackle this issue, we suggest representing a dense orthogonal matrix as a composition of several sparse orthogonal matrices. For a unified and conceptually accessible framework to analyze the sparsity structures in orthogonal matrix decompositions, we recast the generation of a dense orthogonal matrix within OFT as an information propagation problem. Specifically, constructing a dense matrix $\bm{R}\in\mathbb{R}^{d\times d}$ by multiplying $m$ square matrices $\bm{R}=\bm{B}_m\bm{B}_{m-1}\cdots\bm{B}_1$ can be conceptualized as sending information across a $d\times (m+1)$ grid, as visualized in Figure~\ref{fig:info_trans}. This analogy arises from recognizing that a dense $d$-dimensional square matrix establishes complete connectivity between two sets of $d$ nodes. For $\bm{R}$, if an entry $R_{ij}$ equals zero, information from node $j$ is unable to reach node $i$; a nonzero $R_{ij}$, conversely, permits this transmission. Consequently, expressing $\bm{R}$ as a sequence of matrices $\bm{B}_m\bm{B}_{m-1}\cdots\bm{B}_1$ can also be interpreted as modeling a multi-stage information exchange process according to the graphs determined by each $\bm{B}_i$. The flow of information is sequential, proceeding through $\bm{B}_1$ first and culminating with $\bm{B}_n$. As shown in Figure~\ref{fig:info_trans}, consider the factorization $\bm{R}=\bm{B}_5\bm{B}_4\bm{B}_3\bm{B}_2\bm{B}_1$, whose sparsity and associated graph are illustrated. The depicted graph arises from unfolding the sequence of matrix multiplications: here, each $\bm{B}_i$ serves as the adjacency matrix linking the $i$-th level nodes to those at level $(i+1)$. In more detail, entry $(j_1,j_2)$ in $\bm{B}_i$ specifies if there is a directed connection from node $j_2$ at level $i$ to node $j_1$ at level $(i+1)$, with zeros indicating the absence of a link. For the product $\bm{B}_5\bm{B}_4\bm{B}_3\bm{B}_2\bm{B}_1$ to yield a dense matrix, information should flow from every node at the first level to all nodes at the sixth level. If we restrict our attention to $\bm{R}=\bm{B}_4\bm{B}_3\bm{B}_2\bm{B}_1$, thus focusing on information passage from the first-level (source) nodes to fifth-level (receiver) nodes, we observe, for example, that node $1$ cannot send information to node $3$, signifying that the $(3,1)$ entry in the resulting sparsity pattern is zero. The same principle linking the induced graph to the sparsity structure applies to $\bm{R}=\bm{B}_3\bm{B}_2\bm{B}_1$ and $\bm{R}=\bm{B}_2\bm{B}_1$. In the general case, for a matrix $\bm{R}\in\mathbb{R}^{d\times d}$ to be fully dense, the $m$ component matrices $\bm{B}_m,\ldots,\bm{B}_1$ must form a network of directed edges within a $d\times (m+1)$ grid, where each directed edge only links nodes between consecutive levels (i.e., adjacent columns), and this structure ensures that all information from the initial nodes propagates to every node at the final level.

Adopting the perspective of information transmission allows for a clearer interpretation of the sparsity patterns characterizing factorization matrices in OFT. As depicted in Figure~\ref{fig:blk_diag}, the matrix $\bm{R}$ in standard OFT possesses a block-diagonal layout. Although this structure achieves parameter reduction, it falls short in realizing a dense matrix $\bm{R}$. The objective, therefore, is to synthesize a dense orthogonal matrix from $m$ sparse orthogonal factors, while minimizing the number of active trainable parameters. Within the framework of information transmission, two principal requirements emerge: \emph{(i)} \textbf{dense connectivity}, wherein every node at the initial layer should be linked—through at least one pathway—to every node at the final layer; and \emph{(ii)} \textbf{minimal free edges}, meaning the total number of trainable edges must be minimized subject to orthogonality constraints. It is important to recognize that enforcing orthogonality imposes intricate restrictions on the edges linking consecutive layers. For instance, each matrix $\bm{B}_i$ must be of full rank—a prerequisite for orthogonality—which entails that there must exist $d$ edges forming a bijection between nodes at the $i$-th and $(i+1)$-th layers. Thus, the edge count between neighboring layers is bounded below by $d$ (\emph{e.g.}, four in the scenario illustrated in Figure~\ref{fig:info_trans}). These $d$ essential connections are dictated by orthogonality and should be excluded from the count of trainable edges, as their corresponding matrix entries are not subject to learning (\emph{e.g.}, for a $d\times d$ orthogonal matrix with $d$ nonzero entries, those entries are necessarily $\pm 1$). Given that orthogonal matrices can be parameterized with fewer values than unrestricted matrices, orthogonality introduces further dependencies among the edge configurations. For example, in the $2\times 2$ case, if the $(1,1)$ entry is zero, the $(2,2)$ entry must also be zero (\emph{i.e.}, removal of one edge can force the removal of another). Accordingly, orthogonality can, depending on the specific pattern, either add or eliminate certain edges. By visualizing the pattern of non-zero elements in the resulting orthogonal matrix, the information transmission approach proves particularly advantageous for OFT, since the focus is solely on which entries of $\bm{R}$ are non-zero and learnable, rather than on their particular values.

A straightforward fully-connected arrangement between two layers requires $\mathcal{O}(d^2)$ connections (corresponding to one dense orthogonal matrix), resulting in $d^2 - d$ edges (for instance, when $d=4$, there are $12$ edges). Figure~\ref{fig:info_trans} depicts a realizable matrix factorization that employs $10$ edges in total—this is fewer than what a dense orthogonal matrix necessitates. This graphical framework allows for an analysis of OFT’s parameter efficiency, facilitating the design of effective factorizations. In constructing efficient inter-level connectivity, we take inspiration from the butterfly graphs~\cite{cooley1965algorithm}, a topological pattern in the Cooley-Tukey algorithm, which enables a complete connection of $d$ sources to $d$ targets with only $\mathcal{O}(d \log d)$ connections. Notably, the arrangement shown in Figure~\ref{fig:info_trans} achieves dense connectivity using $10$ edges, but a butterfly configuration manages this with merely $8$. We proceed by describing the butterfly structure’s advantage in terms of parameter reduction.

\section{Orthogonal Parameterization by Butterfly Factorization}

The butterfly structure originated from the Cooley-Tukey algorithm as a means to execute the fast Fourier transform. In this context, a local perturbation in the frequency domain propagates as a global shift in the spatial domain, resembling our challenge of information transmission, where any node in the initial layer can relay information to all nodes at the final layer. Furthermore, butterfly structures are well-established as efficient topologies in computer network design~\cite{solihin2015fundamentals,li2011linear}. Setting $k\geq 2$ and restricting $k$ to be a power of $2$, we introduce the \emph{butterfly factor} $\bm{B}^F(k)$ by

\begin{equation}\label{eq:but_factor}
\begin{aligned}
    \bm{B}^F(k)=\begin{bmatrix}
    \text{diag}(\bm{d}_1) & \text{diag}(\bm{d}_2) \\
    \text{diag}(\bm{d}_3) & \text{diag}(\bm{d}_4)
    \end{bmatrix}\in\mathbb{R}^{k\times k},
\end{aligned}
\end{equation}

where each $\bm{d}_i\in\mathbb{R}^{\frac{k}{2}}$ for $i=1,2,3,4$ denotes a vector. For $d=2^N$, the $d$-dimensional \emph{butterfly matrix} $\bm{B}(d)\in\mathbb{R}^{d\times d}$ is defined recursively as follows:

\begin{equation}
\begin{aligned}
    \!\!\bm{B}(d)=\tilde{\bm{B}}(d,d)\!\cdot\!\begin{bmatrix}
    \bm{B}_1(\frac{d}{2})\!\!\!\!\! & \bm{0} \\
    \bm{0}\!\!\!\!\! &\bm{B}_2(\frac{d}{2})
    \end{bmatrix}= \tilde{\bm{B}}(d,d)\tilde{\bm{B}}(d,\frac{d}{2})\cdots\tilde{\bm{B}}(d,2),
\end{aligned}
\end{equation}

Here, $\bm{B}_1(\frac{d}{2})$ and $\bm{B}_2(\frac{d}{2})$ denote two butterfly matrices of dimension $\frac{d}{2}$. We introduce the \emph{butterfly component} as $\thickmuskip=2mu \medmuskip=2mu \tilde{\bm{B}}(d,k)=\text{diag}(\bm{B}^F_1(k),\cdots,\bm{B}^F_{d/k}(k))$, which constitutes a block-diagonal $d\times d$ matrix with each block of size $k$. Here, each $\bm{B}^F_i(k)$ refers to the butterfly factors previously described in Equation~\ref{eq:but_factor}. This structure allows us to employ the butterfly matrix as a way to parameterize an orthogonal matrix: it suffices to ensure that each multiplicative component $\tilde{\bm{B}}(d,k)$ within $\bm{B}(d)$ is itself orthogonal. To explore this, we start with the specific case where the block size is $2$, i.e., $\tilde{\bm{B}}(d,2)$, and observe that we can guarantee its orthogonality by using either the Cayley transform or $2$-dimensional rotations for each block—consistent with techniques in \cite{liu2021orthogonal,qiu2023controlling}. Notably, every butterfly component's sparsity pattern corresponds to some permutation of the non-zero structure in $\tilde{\bm{B}}(d,2)$, so each can also be readily constructed as an orthogonal matrix. Consequently, this approach yields a highly efficient orthogonal matrix parameterization built from numerous $2\times 2$ orthogonal blocks. Extending the framework as in \cite{chen2022pixelated}, we define a block butterfly component $\tilde{\bm{B}}^b(d,k)$, in which each element of $\bm{d}_i$ becomes a $b\times b$ submatrix. To maintain the orthogonality of $\tilde{\bm{B}}^b(d,2)$, we parameterize each $2b\times 2b$ block as an orthogonal matrix. The sparsity pattern of other butterfly components $\tilde{\bm{B}}^b(d,k)$ with $k>2$ is realized by block-wise permutations of the pattern in $\tilde{\bm{B}}^b(d,2)$, allowing for the same parameterization strategy to enforce orthogonality. Collectively, these elements form the foundation for the BOFT forward computation:

\begin{equation}
    \begin{aligned}\nonumber
    \bm{z}=\big(\bm{R}(m,b)\cdot\bm{W}^0\big)^\top\bm{x},~~~~\text{s.t.}~\bigg{\{}\bm{R}(m,b)=\prod_{i=1}^m \tilde{\bm{B}}^b_{(d,i)}~~\&~~\underbrace{\big(\tilde{\bm{B}}^b_{(d,j)}\big)^\top\tilde{\bm{B}}^b_{(d,j)}=\tilde{\bm{B}}^b_{(d,j)}\big(\tilde{\bm{B}}^b_{(d,j)}\big)^\top\!=\bm{I}_d}_{\forall j\in [1, m]}\bigg{\}},
    \end{aligned}
\end{equation}

Here, for brevity, $\tilde{\bm{B}}^b(d,2^{m - i+1})$ is written as $\tilde{\bm{B}}^b_{(d,i)}$, and $\bm{I}_d$ denotes the $d \times d$ identity matrix. The orthogonal matrix $\bm{R}(m,b)\in\mathbb{R}^{d\times d}$ is constructed as the product of $m$ orthogonal butterfly matrices. For clarity, we denote BOFT employing $\bm{R}(m,\frac{b}{2})$ as BOFT($m$,$b$), given $b\geq 2$. In the special case when $\thickmuskip=2mu \medmuskip=2mu m=1$, BOFT($1$,$b$) is equivalent to the block-diagonal OFT~\cite{qiu2023controlling} having block size $b$. Setting BOFT($1$,$d$) retrieves the original OFT~\cite{qiu2023controlling} with an unconstrained orthogonal matrix. When $m=\log \frac{2d}{b}$, BOFT($\log \frac{2d}{b}$,$b$) incorporates the block butterfly matrix $\bm{B}^{\frac{b}{2}}(d)$ for $\bm{R}$, resulting in a dense orthogonal matrix. Typically, BOFT($m$,$b$) introduces $\frac{1}{2}(b-1)dm$ trainable parameters when tuning a linear layer of dimension $d\times n$. For the pure butterfly structure ($m=\log d,~b=2$), the parameter count is $\mathcal{O}(d\log d)$. On the other hand, the full dense OFT utilizes $\mathcal{O}(d^2)$ parameters, while the block-diagonal variant with $r$ blocks employs $\mathcal{O}(bd)$. Notably, achieving a dense orthogonal matrix in standard OFT requires choosing $b=d$, whereas BOFT can produce dense representations regardless of the value of $b$.

\textbf{Identity Initialization in BOFT}. Fine-tuning methods frequently initialize from the parameters of a pretrained network to ensure minimal deviation. LoRA, for instance, initializes its low-rank updates as zero matrices. In the case of BOFT, each butterfly block is initially set as the identity matrix, meaning the corresponding skew-symmetric matrix in the Cayley transform starts with zeros.

\textbf{BOFT-specific Multiplicative Dropout}. To enhance robustness against overfitting, LoRA~\cite{hulora2022} integrates Dropout on its low-rank updates, leveraging the additive nature of the updates. The standard Dropout~\cite{srivastava2014dropout} is not directly compatible with BOFT, since BOFT updates are multiplicative. Therefore, we introduce a multiplicative Dropout mechanism tailored to BOFT. Specifically, since the orthogonal $\bm{R}(m,b)$ is assembled from $m$ butterfly matrices, each being permutable into block-diagonal form with blocks of size $2b\times 2b$, we randomly select $p_1$ percent of the butterfly blocks and $p_2$ percent of blocks within each butterfly; these selected blocks are then substituted with identity matrices to perform multiplicative Dropout.

\section{Intriguing Insights and Discussions}

\textbf{Expressivity of BOFT}. The butterfly structure, together with permutation operations, is capable of exactly representing several well-known fast linear transforms~\cite{dao2019learning,dao2019kaleidoscope}, such as the fast Fourier transform and the Hadamard transform. However, the extent to which our orthogonal butterfly matrix can approximate an arbitrary orthogonal matrix is still undetermined. To investigate this, we perform a simulation targeting the approximation of a randomly generated dense orthogonal matrix~\cite{anderson1987generation} of dimension $512\times 512$. In Figure~\ref{fig:para_eff}, outcomes are averaged across 10 different random seeds. The y-axis quantifies the approximation error, while the x-axis represents the number of effective trainable parameters. Each same-colored curve corresponds to BOFT instances with an identical block size; the curve’s leftmost marker denotes BOFT($1$,$b$) (i.e., the baseline block-diagonal OFT with block size $b$). Generally, BOFT demonstrates superior parameter efficiency over OFT. For instance, BOFT($9$,$2$) achieves greater expressivity than BOFT($1$,$16$) while employing fewer parameters. More broadly, configurations of BOFT with smaller $b$ and larger $m$ tend to achieve higher parameter efficiency; BOFT($6$,$4$), for example, utilizes significantly fewer parameters yet attains an approximation error comparable to that of BOFT($2$,$16$). These results suggest that butterfly matrices define a more structured subset within the space of orthogonal matrices compared to the block-diagonal approach, providing BOFT with provably enhanced expressivity relative to OFT with equivalent block size.

\begin{theorem}[Expressivity of BOFT]\label{thm:exp_boft}
    BOFT achieves greater expressivity than OFT for an identical block size. To represent every orthogonal matrix of dimension $d$, it suffices to take products of butterfly matrices as $\bm{B}_{d-1,1}(d)\bm{B}^\top_{d-1,2}(d)\cdots\bm{B}_{1,1}(d)\bm{B}^\top_{1,2}(d)$, where each $\bm{B}_{i,j}(d),\forall i,\forall j$ is itself a butterfly matrix.
\end{theorem}

Theorem~\ref{thm:exp_boft} provides a natural way to generalize BOFT: the resultant orthogonal matrix can be expressed as $\thickmuskip=2mu \medmuskip=2mu \bm{R}^G(m_1,b_1,m_2,b_2,l)=\bm{R}_{l,1}(m_1,b_1)\bm{R}_{l,2}^T(m_2,b_2)\cdots\bm{R}_{1,1}(m_1,b_1)\bm{R}_{1,2}^T(m_2,b_2)$, where $\bm{R}_{i,j}^T(m,b)$ denotes the orthogonal matrices employed in BOFT. In particular, when $m_1=m_2=\log d$, $b_1=b_2=2$, and $l=d-1$, $\bm{R}^G(m_1,b_1,m_2,b_2,l)$ spans the entire orthogonal group. This form of matrix composition has also been referred to as the kaleidoscope hierarchy~\cite{dao2019kaleidoscope}. Nevertheless, greater expressivity does not automatically guarantee improved performance during finetuning—indeed, full finetuning, despite covering the complete expressive range, may yield suboptimal outcomes. Thus, achieving a balance between expressivity and imposed regularity becomes central to enhancing model generalization in finetuning scenarios. BOFT expands the parameter search space for finetuning by introducing structural priors, allowing for a more favorable compromise between expressivity and regularity.

\textbf{Spectral properties}. Compared to LoRA, orthogonal finetuning tends to achieve superior spectral characteristics due to its ability to exactly retain the spectral norm of the original pretrained weight matrix $\bm{W}^0$. To see why, consider the singular value decomposition $\bm{W}^0=\bm{U}\bm{\Sigma}\bm{V}^\top$, where $\bm{U}$ and $\bm{V}$ are orthogonal matrices and $\bm{\Sigma}$ is a diagonal matrix of singular values. Both OFT and BOFT implement an orthogonal transformation by left-multiplying with a matrix $\bm{R}$, yielding new weights of the form $\bm{R}\bm{U}\bm{\Sigma}\bm{V}^\top$. This operation leaves the dominant singular value (i.e., the spectral norm of $\bm{W}^0$) unchanged. Such exact preservation is recognized as beneficial for improving training stability and promoting generalization~\cite{miyato2018spectral,yoshida2017spectral}. Additional mathematical properties related to this are provided in Appendix~\ref{app:sn_property}.

\textbf{Orthogonal finetuning as learning bilinear similarity}. The BOFT framework may also be interpreted as learning a bilinear similarity, specifically in the form $\bm{w}^0_i\bm{R}\bm{x}$, where $\bm{w}^0_i$ represents the $i$-th neuron (that is, column vector) of $\bm{W}^0$. In this context, BOFT focuses on learning the matrix $\bm{R}$, which defines the bilinear similarity under the strict constraint that $\bm{R}$ remains orthogonal. This links BOFT closely with the principles of distance metric learning~\cite{xing2002distance} and the theory of bilinear forms~\cite{roman2005advanced}.

\textbf{Inductive bias and generalization in BOFT}. The $\bm{R}(m,b)$ matrix used in BOFT often comes from a specifically structured subset within the orthogonal group, thereby limiting the function space BOFT can explore and introducing a corresponding inductive bias. The structured patterns enforced by butterfly factorization, which align with patterns found in numerous canonical linear transforms such as the discrete Fourier transform, discrete sine and cosine transforms, and Hadamard transform~\cite{dao2019kaleidoscope}, are argued to foster improved generalization. Further, the use of sparse matrix factorization in BOFT can impart additional implicit forms of inductive bias~\cite{gunasekar2017implicit,li2018algorithmic,liu2019neural}.

\textbf{Comparison to butterfly-based sparse training}. Several prior studies~\cite{dao2019learning,dao2019kaleidoscope,chen2022pixelated,dao2022monarch} have explored sparse training methodologies utilizing butterfly parameterization. These approaches predominantly reparameterize weight matrices directly through the butterfly structure and initiate the neural network training process anew. In the work by \cite{dao2022monarch}, pretrained weights are refined by initially projecting them onto a modified form of butterfly matrices, subsequently optimizing these projected components for specific downstream tasks. In contrast, BOFT introduces a distinct finetuning mechanism, applying transformations to the weights by means of layer-shared weight matrices.

\input{rebuttal-results}

\section{Concluding Remarks and Limitations}

In this work, we introduce Orthogonal Butterfly, a general method for parameter-efficient finetuning of foundation models grounded in the butterfly structure. Our central observation for enhancing parameter-efficiency lies in representing a dense orthogonal matrix as a product of several sparse orthogonal matrices. To facilitate practical matrix factorizations, we develop a graphical information transmission framework. Within this framework, we identify the butterfly structure as a highly effective means for sparse orthogonal matrix decomposition, aligning with our aims. Our empirical results indicate that BOFT yields strong finetuning performance on large language models, vision models, and text-to-image generative architectures. These findings underscore the capability of BOFT as a broadly applicable finetuning technique.

Nevertheless, there are certain limitations to BOFT. Notably, constructing the final orthogonal matrix as a product of multiple orthogonal matrices incurs a modest increase in training runtime overhead compared to OFT. Optimizing the training efficiency of BOFT is an open area for further research. Fortunately, after finetuning is complete, the orthogonal matrices generated by BOFT can be consolidated within the pretrained model with no added inference time. Additionally, the question remains as to whether the butterfly topology offers the optimal strategy for information transmission. Our information transmission framework also encourages exploration inspired by the field of computer networking, where analysis of network topology efficiency is a core concern. This connection opens the door to designing alternative, potentially more efficient network structures for constructing dense orthogonal matrices.

\end{document}